\begingroup
\color{borrador2}
\textcolor{borrador2}{La siguiente figura amplía como alternativa la vista del servidor para mostrar sus dependencias y los eventos del flujo de respuesta.}

\begin{figure}[H]
  \centering
  \resizebox{0.98\textwidth}{!}{%
    \begin{tikzpicture}[
      nodo/.style={draw=borrador2, thick, rounded corners=3pt, fill=white,
                   text=borrador2, align=center, font=\small,
                   minimum width=3.4cm, minimum height=1.0cm},
      modulo/.style={nodo, fill=borrador2!7, minimum width=4.2cm,
                     font=\footnotesize},
      externo/.style={nodo, densely dashed},
      dato/.style={nodo, fill=borrador2!10, font=\footnotesize\ttfamily},
      grupo/.style={draw=borrador2, very thick, rounded corners=5pt, inner sep=9pt},
      titulo/.style={font=\small\bfseries, text=borrador2},
      flecha/.style={-{Latex[length=2.2mm]}, thick, draw=borrador2},
      doble/.style={{Latex[length=2.2mm]}-{Latex[length=2.2mm]}, thick, draw=borrador2},
    ]
      \node[nodo, minimum width=4.5cm, minimum height=1.4cm] (browser) at (0,0)
        {Navegador\\{\scriptsize HTML · CSS · JavaScript}};

      \node[modulo] (http) at (6.0,2.7) {Rutas HTTP y archivos estáticos};
      \node[modulo] (conv) at (6.0,0.9) {Conversación y decisión de ámbito};
      \node[modulo] (rag) at (6.0,-0.9) {Recuperación, generación y verificación};
      \node[modulo] (aux) at (6.0,-2.7) {Sugerencias y registro de uso};

      \node[externo] (ollama) at (12.5,1.4) {Ollama\\{\scriptsize proceso de inferencia local}};
      \node[externo, cylinder, shape border rotate=90, aspect=0.18] (indice) at (12.5,-0.8)
        {LanceDB};
      \node[dato] (jsonl) at (12.5,-3.0) {registro\_chat.jsonl};

      \draw[doble] (browser) -- node[above, font=\scriptsize, text=borrador2, align=center]
        {\texttt{/api/preguntar}\\fuentes · partes · borrar · sugerencias} (http);
      \draw[flecha] (http) -- (conv);
      \draw[flecha] (conv) -- (rag);
      \draw[flecha] (rag) -- (aux);
      \draw[doble] (conv.east) -- (ollama.west);
      \draw[doble] (rag.east) -- (ollama.west);
      \draw[doble] (rag.east) -- (indice.west);
      \draw[flecha] (aux.east) -- (jsonl.west);

      \begin{scope}[on background layer]
        \node[grupo, fit=(http)(conv)(rag)(aux),
              label={[titulo]north:Único proceso de servidor Python}] {};
      \end{scope}

      \node[font=\scriptsize, text=borrador2, align=center, anchor=north]
        at (0,-1.0)
        {La interfaz no requiere cuenta;\\el estado de conversación vive en el servidor.};
    \end{tikzpicture}%
  }
  \caption[Alternativa de la arquitectura web]{\textcolor{borrador2}{Versión alternativa de la arquitectura web. Un proceso Python sirve la interfaz y coordina conversación, ámbito, RAG, sugerencias y registro; LanceDB, Ollama y el fichero JSONL se muestran como dependencias externas al proceso. Fuente: Elaboración propia.}}
  \label{fig:arquitectura_web_alt}
\end{figure}
\endgroup
